\documentclass[11pt,a4paper]{article}

%====================================================
% PACKAGES
%====================================================

\usepackage[utf8]{inputenc}
\usepackage[T1]{fontenc}
\usepackage[french]{babel}

\usepackage{geometry}
\geometry{margin=2.5cm}

\usepackage{amsmath,amssymb,amsthm}
\usepackage{mathtools}

\usepackage{graphicx}
\usepackage{float}

\usepackage{booktabs}
\usepackage{multirow}
\usepackage{array}
\usepackage{longtable}

\usepackage{enumitem}

\usepackage[dvipsnames]{xcolor}

\usepackage{caption}
\usepackage{subcaption}

\usepackage{hyperref}
\hypersetup{
    colorlinks=true,
    linkcolor=Black,
    urlcolor=Black,
    citecolor=Black
}

\usepackage{titlesec}
\usepackage{fancyhdr}
\usepackage{tcolorbox}

\usepackage{algorithm}
\usepackage{algpseudocode}

\usepackage{setspace}
\usepackage{tikz}
\usepackage{pgfplots}

%====================================================
% STYLE
%====================================================

\setlength{\parindent}{0pt}
\setlength{\parskip}{0.7em}

\pagestyle{fancy}
\fancyhf{}
\fancyhead[L]{Orienteering Problem}
\fancyhead[R]{}
\fancyfoot[C]{\thepage}

\titleformat{\section}
{\Large\bfseries\color{Black}}
{\thesection.}{0.5em}{}

\titleformat{\subsection}
{\large\bfseries\color{Black}}
{\thesubsection.}{0.5em}{}

\titleformat{\subsubsection}
{\normalsize\bfseries\color{Black}}
{\thesubsubsection.}{0.5em}{}

\newcommand{\figplaceholder}[1]{%
\fbox{%
\parbox{0.92\linewidth}{%
\centering
\vspace{1.8cm}
\textbf{#1}
\vspace{1.8cm}
}%
}%
}

%====================================================
% TITLE
%====================================================

\title{
\vspace{-1cm}
\textbf{\Huge Étude du Problème du Voyageur à Profits}\\
\vspace{0.3cm}
\Large Orienteering Problem (OP)
}

\author{
\textbf{Ela SARHANI}
}

\date{\today}

%====================================================
% DOCUMENT
%====================================================

\begin{document}

\begin{titlepage}

\begin{center}

\vspace*{2.5cm}

%------------------------------------------------
% TITLE
%------------------------------------------------

{\Huge\bfseries Étude du Problème du Voyageur à Profits \par}

\vspace{0.6cm}

{\LARGE\itshape Orienteering Problem (OP)\par}

\vspace{1.8cm}

\hrule
\vspace{0.5cm}
{\Large Rapport de Projet}
\vspace{0.5cm}
\hrule

\vspace{3cm}

%------------------------------------------------
% INFORMATION TABLE
%------------------------------------------------

\renewcommand{\arraystretch}{1.6}

\begin{tabular}{>{\bfseries}r l}
Réalisé par : & Ela SARHANI \\

Matière : & Optimisation Combinatoire \\

Année universitaire : & 2025 -- 2026 \\
\end{tabular}

\vfill

%------------------------------------------------
% FOOTER
%------------------------------------------------

{\large \today\par}

\vspace{1cm}

\end{center}

\end{titlepage}

%====================================================
% ABSTRACT
%====================================================

\vspace*{\fill}

\begin{center}
\begin{minipage}{1\textwidth}

\begin{abstract}

Ce rapport présente une étude complète du \emph{Orienteering Problem} (OP), problème de routage avec profits soumis à une contrainte de budget. L'objectif est de sélectionner et d'ordonner un sous-ensemble de sommets afin de maximiser le profit total tout en respectant une borne de coût.

Le travail combine une modélisation mathématique rigoureuse, trois méthodes exactes distinctes, trois métaheuristiques adaptées, puis une comparaison expérimentale sur plusieurs instances. L'ensemble est analysé dans une perspective scientifique, en mettant l'accent sur la structure combinatoire du problème, les choix algorithmiques et l'interprétation critique des résultats.

Une attention particulière est également portée à la visualisation des solutions et à l'analyse comparative des comportements algorithmiques afin de mieux comprendre les mécanismes internes de recherche.

\end{abstract}

\end{minipage}
\end{center}

\vspace*{\fill}

\vspace{0.4cm}

\newpage

\tableofcontents

\newpage

%====================================================
% INTRODUCTION
%====================================================

\section{Introduction et choix du problème}

Le problème choisi est le \emph{Orienteering Problem} (OP), souvent appelé le problème du \emph{voyageur à profits} ou \emph{Prize Collecting Traveling Salesman Problem}. Ce choix est fondamental pour une étude comparative approfondie comparant méthodes exactes et métaheuristiques, car le problème est à la fois suffisamment riche pour nécessiter une vraie modélisation combinatoire sophistiquée et suffisamment structuré pour permettre des approches algorithmiques variées et complémentaires.

Contrairement au célèbre \emph{Traveling Salesman Problem} (TSP), l'OP n'impose pas de visiter tous les sommets du graphe. Il introduit ainsi une dimension décisionnelle supplémentaire : non seulement il faut déterminer l'ordre de visite des sommets sélectionnés, mais il faut aussi décider \emph{quels} sommets méritent vraiment d'être inclus dans la tournée. Cette complexité additionnelle confère au problème un statut de problème \emph{NP-difficile} au sens fort.

L'OP combine deux décisions essentielles et souvent contradictoires :

\begin{itemize}
    \item \textbf{la sélection des sommets à visiter} : chaque sommet offre un profit, mais ajouter une visite augmente le coût de déplacement ;
    \item \textbf{l'ordonnancement optimal de ces sommets} : une fois les sommets sélectionnés, il faut trouver la meilleure séquence de visite.
\end{itemize}

Cette double dimension crée un compromis direct et inévitable entre profit collecté et coût de déplacement. Le décideur doit donc arbitrer : collecter beaucoup de profits mais en voyageant davantage, ou se limiter à quelques sommets à haut profit proches du dépôt. Ce conflit rend le problème représentatif des difficultés classiques et les plus profondes de l'optimisation combinatoire moderne.

\subsection{Applicabilités concrètes}

L'OP modélise naturellement une large gamme de situations réelles :

\begin{itemize}
    \item \textbf{Tourisme et loisirs} : planifier une visite complète de monuments et musées dans une région, en maximisant les visites intéressantes tout en restant dans une limite de temps ou de budget ;
    \item \textbf{Surveillance et inspection} : programmer des missions de surveillance autonomes sur un territoire, en sélectionnant les zones prioritaires à couvrir ;
    \item \textbf{Logistique et collecte} : organiser des tournées de collecte sélective, en décidant quels clients servir en fonction de leur profitabilité et de leur localisation ;
    \item \textbf{Systèmes énergétiques} : robots ou véhicules autonomes avec batterie limitée, cherchant à maximiser les tâches accomplies avant épuisement énergétique ;
    \item \textbf{Planification de missions spatiales ou scientifiques} : sélectionner les points d'intérêt à explorer avec une contrainte de carburant ou de temps.
\end{itemize}

\subsection{Structure combinatoire et enjeux de modélisation}

La difficulté principale de l'OP vient du fait que l'espace des solutions croît très rapidement avec le nombre de sommets. Chaque solution admissible correspond à une tournée partant d'un dépôt, visitant un sous-ensemble de sommets, puis revenant au dépôt.

Le nombre de possibilités combine :

\begin{itemize}
    \item le choix des sommets visités : $2^{n-1}$ possibilités ;
    \item l'ordre de visite de ces sommets : jusqu'à $(n-1)!$ ordres différents.
\end{itemize}

Le problème appartient donc à la famille des problèmes combinatoires difficiles, avec un espace de solutions de taille $O(2^n \cdot n!)$ dans le pire cas.

Du point de vue de la modélisation, plusieurs contraintes doivent être simultanément satisfaites :

\begin{itemize}
    \item cohérence de la tournée (chaque sommet visité une seule fois) ;
    \item respect du budget global ;
    \item unicité des visites ;
    \item élimination des sous-tours disjoints.
\end{itemize}

\subsection{Pourquoi le problème se prête aux approches exactes et métaheuristiques}

\subsubsection*{Adaptation technique aux approches exactes}

L'OP se prête remarquablement bien aux approches exactes pour plusieurs raisons techniques profondes :

\begin{enumerate}
    \item \textbf{Structure d'arborescence naturelle} : Le problème peut être décomposé naturellement par exploration arborescente. Chaque branche de l'arbre représente une décision : inclure ou exclure un sommet, ou changer l'ordre de visite. Cette structure permet d'appliquer élégamment le \emph{Branch and Bound}.
    
    \item \textbf{Construction de bornes efficaces} : Il est possible de calculer rapidement des bornes supérieures utiles. Par exemple, assumer que tous les sommets restants peuvent être visités au coût minimum, ou utiliser une relaxation linéaire. Ces bornes permettent de couper de larges portions de l'arbre de recherche.
    
    \item \textbf{Structure de programmation linéaire entière bien posée} : L'OP peut se formuler clairement en variables binaires avec des contraintes linéaires (conservation de flot, élimination de sous-tours). Cette formulation permet l'usage de solveurs de programmation linéaire entière et de techniques de coupes avancées.
    
    \item \textbf{Programmation dynamique applicable} : L'espace d'état (position courante, ensemble de sommets visités, budget consommé) est fini et peut être mémorisé. La relation de récurrence entre états est claire, ce qui rend la DP directement applicable sur des instances modérées.
    
    \item \textbf{Complexité intermédiaire} : Contrairement au TSP classique où il faut visiter tous les sommets, l'OP offre une flexibilité qui, sur petites instances, rend les méthodes exactes tractables. L'espace de recherche reste gigantesque mais les bonnes bornes et coupes réduisent drastiquement ce dernier en pratique.
\end{enumerate}

\subsubsection*{Adaptation technique aux métaheuristiques}

L'OP est également particulièrement bien adapté aux métaheuristiques pour des raisons structurelles et pratiques :

\begin{enumerate}
    \item \textbf{Évaluation rapide} : Évaluer une solution (vérifier la faisabilité, calculer le profit total et le coût) se fait en temps $O(n)$. Cela permet aux métaheuristiques d'explorer rapidement de nombreuses solutions candidates.
    
    \item \textbf{Voisinages naturels et riches} : Il existe plusieurs définitions naturelles de voisinage :
    \begin{itemize}
        \item Insertion/suppression d'un sommet ;
        \item Inversion de deux sommets (2-opt) ;
        \item Remplacement d'un sommet par un autre ;
        \item Permutation locale de la séquence.
    \end{itemize}
    Ces voisinages sont simples à implémenter mais offrent suffisamment de diversité pour explorer l'espace des solutions.
    
    \item \textbf{Gestion facile de la faisabilité} : Contrairement aux problèmes avec de nombreuses contraintes complexes, la faisabilité dans l'OP repose principalement sur une seule contrainte : le budget total. Cela facilite la construction de solutions initialisées (une complètion gloutonne produit automatiquement une solution faisable) et les réparations de solutions violant la contrainte.
    
    \item \textbf{Existence d'améliorations locales fréquentes} : Le paysage de fitness de l'OP contient beaucoup de vallées et de pics locaux. Les métaheuristiques qui explorent ce paysage trouvent facilement des améliorations, ce qui rend la recherche « guidée » et productive.
    
    \item \textbf{Passage à l'échelle} : Bien que le problème soit NP-difficile, la structure permet aux métaheuristiques de maintenir des performances raisonnables même sur des instances contenant des centaines de sommets.
\end{enumerate}

\textbf{Conclusion} : L'OP représente un cas d'étude idéal : il est assez riche pour exercer toutes les techniques, assez structuré pour que chaque approche soit efficace, et assez difficile pour que le contraste entre les familles de méthodes soit dramatique et instructif.

\begin{figure}[H]
\centering
\figplaceholder{Capture du hub de visualisation global : outputs/index.html avec les deux menus déroulants et chargement des viewers}
\caption{Vue d'ensemble de l'interface de visualisation utilisée pour comparer les méthodes.}
\label{fig:hub}
\end{figure}

%====================================================
% FORMAL DEFINITION
%====================================================

\section{Définition formelle du problème}

Soit un graphe complet orienté :

\[
G=(V,E)
\]

où :

\begin{itemize}
    \item $V=\{0,1,\dots,n\}$ représente l'ensemble des sommets ;
    \item $0$ représente le dépôt ;
    \item chaque sommet $i$ possède un profit $p_i>0$ ;
    \item chaque arc $(i,j)$ possède un coût $c_{ij}\ge0$.
\end{itemize}

Le budget maximal autorisé est noté $B$.

L'objectif consiste à construire une tournée :

\begin{itemize}
    \item partant du dépôt ;
    \item visitant un sous-ensemble de sommets ;
    \item revenant au dépôt ;
    \item respectant le budget ;
    \item maximisant le profit total collecté.
\end{itemize}

\subsection{Ensembles, paramètres et variables}

\begin{table}[H]
\centering
\caption{Notation mathématique utilisée}
\begin{tabular}{ll}
\toprule
\textbf{Symbole} & \textbf{Description} \\
\midrule
$V$ & Ensemble des sommets \\
$E$ & Ensemble des arcs \\
$0$ & Dépôt de départ et d'arrivée \\
$p_i$ & Profit associé au sommet $i$ \\
$c_{ij}$ & Coût de déplacement entre $i$ et $j$ \\
$B$ & Budget maximal autorisé \\
$x_{ij}$ & Variable binaire d'utilisation d'arc \\
$y_i$ & Variable binaire de visite \\
$u_i$ & Variable MTZ d'ordre \\
\bottomrule
\end{tabular}
\end{table}

\subsection{Fonction objectif}

La fonction objectif consiste à maximiser le profit total collecté :

\[
\boxed{
\max \left(
\sum_{i\in V\setminus\{0\}} p_i y_i
\right)
}
\]

où :

\begin{itemize}
    \item $p_i$ représente le profit du sommet $i$ ;
    \item $y_i=1$ si le sommet est visité ;
    \item $y_i=0$ sinon.
\end{itemize}

\subsection{Contraintes}

Une formulation classique de l'OP s'écrit :

\[
\sum_{j\in V\setminus\{0\}} x_{0j}=1
\]

\[
\sum_{i\in V\setminus\{0\}} x_{i0}=1
\]

\[
\sum_{j\in V,\, j\neq i}x_{ij}=y_i
\qquad
\forall i\in V\setminus\{0\}
\]

\[
\sum_{j\in V,\, j\neq i}x_{ji}=y_i
\qquad
\forall i\in V\setminus\{0\}
\]

\[
\sum_{i\in V}\sum_{j\in V,\, j\neq i}
c_{ij}x_{ij}\le B
\]

\[
u_i-u_j+n x_{ij}\le n-1
\]

\[
x_{ij}\in\{0,1\}
\qquad
y_i\in\{0,1\}
\]

Les contraintes MTZ permettent d'éliminer les sous-tours disjoints.

%====================================================
% EXACT METHODS
%====================================================

\section{Méthodes exactes}

Les méthodes exactes cherchent à garantir l'obtention de la solution optimale. Trois approches distinctes ont été étudiées :

\begin{itemize}
    \item Branch and Bound ;
    \item Branch and Cut ;
    \item Programmation dynamique.
\end{itemize}

%====================================================
% BNB
%====================================================

\subsection{Branch and Bound}

\subsubsection*{Justification du choix}

Branch and Bound est choisi comme première méthode exacte car c'est l'algorithme fondamental et le plus direct pour les problèmes d'optimisation combinatoire avec un espace d'énumération naturel. Pour l'OP, il offre un bon compromis entre clarté conceptuelle et efficacité pratique. C'est également l'approche historique de référence pour ce type de problème.

\subsubsection*{Avantages pour ce cas d'usage}

\begin{itemize}
    \item \textbf{Adaptabilité} : Peut être modifié facilement pour ajouter des bornes spécialisées ou des stratégies de branchement adaptées au problème.
    \item \textbf{Efficacité mémoire} : Utilise une exploration en profondeur, ce qui limite la mémoire requise.
    \item \textbf{Bornes efficaces} : Peut calculer rapidement une borne optimiste (profit résiduel maximum) et une borne pessimiste (greedy completion).
    \item \textbf{Visualisation naturelle} : L'arbre de recherche est facile à tracer et visualiser.
\end{itemize}

\subsubsection*{Logique de recherche de la solution optimale}

Branch and Bound fonctionne par exploration arborescente implicite :

1. On maintient une meilleure solution connue (incumbent).
2. À chaque nœud de l'arbre, on calcule une borne supérieure sur le profit maximum atteignable.
3. Si cette borne ne peut pas battre l'incumbent, on élague la branche.
4. Sinon, on branche sur un nouveau sommet à décider (ajouter ou ne pas ajouter à la tournée).

L'algorithme explore exactement les états où la borne reste prometteuse, garantissant ainsi de ne pas manquer la solution optimale.

\subsubsection*{Principe général (pseudocode)}

\begin{algorithm}[H]
\caption{Branch and Bound pour l'OP}
\begin{algorithmic}[1]
\State Initialiser incumbent à -∞
\State Calculer solution greedy, mettre à jour incumbent
\State Créer pile avec le nœud racine (dépôt, aucun sommet visité, budget intact)
\While{pile non vide}
    \State Dépiler un nœud (route partielle, profit courant, coût courant)
    \If{coût courant > budget}
        \State \textbf{continue} \Comment{nœud infaisable}
    \EndIf
    \State Calculer borne supérieure = profit courant + profit résiduel max possible
    \If{borne $\le$ incumbent}
        \State \textbf{continue} \Comment{élague : ne peut pas améliorer}
    \EndIf
    \If{route fermable (coût + retour au dépôt $\le$ budget)}
        \State Mettre à jour incumbent si profit courant > incumbent
    \EndIf
    \For{chaque sommet non visité, ordonné par profit décroissant}
        \State Calculer coût ajout de ce sommet
        \If{coût ajout + coût courant $\le$ budget}
            \State Empiler nouveau nœud (route + sommet, profit + profit\_sommet, coût + coût\_ajout)
        \EndIf
    \EndFor
\EndWhile
\State \textbf{Retourner} incumbent
\end{algorithmic}
\end{algorithm}

\subsubsection*{Complexité}

\begin{itemize}
    \item \textbf{Complexité théorique} : Dans le pire cas, $O(2^n \cdot n!)$ puisqu'on énumère tous les sous-ensembles et tous les ordres.
    \item \textbf{Complexité pratique} : Avec de bonnes bornes, le nombre de nœuds explorés est typiquement très inférieur à $2^n$. Sur des instances de taille modérée (n ≤ 20), l'algorithme est très efficace.
    \item \textbf{Analyse selon la taille} :
    \begin{itemize}
        \item Petites instances (n $\le$ 10) : instantané.
        \item Instances moyennes (n ≈ 15–20) : quelques millisecondes à secondes.
        \item Grandes instances (n > 30) : temps de calcul peut exploser exponentiellement.
    \end{itemize}
\end{itemize}

%====================================================
% B&C
%====================================================

\subsection{Branch and Cut}

\subsubsection*{Justification du choix}

Branch and Cut est choisi comme deuxième approche exacte pour montrer comment renforcer Branch and Bound par l'ajout dynamique de contraintes (coupes). Cette méthode représente l'état de l'art en optimisation entière et démontre l'impact des techniques de renforcement des relaxations.

\subsubsection*{Avantages pour ce cas d'usage}

\begin{itemize}
    \item \textbf{Bornes renforcées} : Les coupes réduisent significativement l'écart entre la solution fractionnaire et la solution entière.
    \item \textbf{Moins de branchement} : Un meilleur contrôle du problème signifie moins de nœuds à explorer.
    \item \textbf{Scalabilité améliorée} : Permet de traiter des instances un peu plus grandes que Branch and Bound seul.
\end{itemize}

\subsubsection*{Logique de recherche}

Branch and Cut alterne entre deux phases :

1. \textbf{Phase de coupe (cutting plane)} : Résoudre la relaxation linéaire. Si la solution est fractionnaire, détecter une contrainte qui coupe cette solution hors-polytope des solutions entières valides. Ajouter cette contrainte et relancer.

2. \textbf{Phase de branchement} : Une fois aucune coupe prometteuse n'existe (ou après un nombre d'itérations), brancher classiquement sur une variable fractionnaire.

Cela renforce progressivement la relaxation, rapprochant la borne inférieure de la solution entière.

\subsubsection*{Principe général (pseudocode)}

\begin{algorithm}[H]
\caption{Branch and Cut pour l'OP}
\begin{algorithmic}[1]
\State Formulation LP initiale : max $\sum p_i y_i$ sous les contraintes linéaires
\State incumbent ← -∞
\State pile ← {nœud root}
\While{pile non vide}
    \State Dépiler nœud
    \State \textbf{repeat}
        \State Résoudre LP du nœud
        \If{LP infaisable}
            \State break
        \EndIf
        \If{LP solution est entière}
            \State Mettre à jour incumbent
            \State break
        \EndIf
        \State Détecter contrainte violée (élimination sous-tour, inégalité de coupe)
        \If{contrainte existe}
            \State Ajouter contrainte à LP
        \Else
            \State break
        \EndIf
    \State \textbf{until} aucune coupe trouvée
    \If{LP borne $\le$ incumbent}
        \State continue
    \EndIf
    \State Brancher sur variable fractionnaire
\EndWhile
\State \textbf{Retourner} incumbent
\end{algorithmic}
\end{algorithm}

\subsubsection*{Complexité}

\begin{itemize}
    \item \textbf{Complexité théorique} : $O(2^n \cdot L \cdot n^3)$ où $L$ est le nombre d'itérations de cutting planes (typiquement petit).
    \item \textbf{Complexité pratique} : Meilleure que Branch and Bound seul grâce aux coupes réduisant l'arborescence. Peut être très rapide sur certaines classes d'instances.
\end{itemize}

%====================================================
% DP
%====================================================

\subsection{Programmation dynamique}

\subsubsection*{Justification du choix}

La Programmation Dynamique est choisie comme troisième approche exacte pour offrir un contraste algorithmique total : contrairement aux deux premières qui sont ascendantes (partent d'une solution partielle et la complètent), la DP est ascendante par définition d'états. Elle montre aussi comment exploiter la structure spécifique du problème (états finis mémorisables).

\subsubsection*{Avantages pour ce cas d'usage}

\begin{itemize}
    \item \textbf{Optimalité garantie} : Si tous les états peuvent être énumérés, la DP trouve l'optimum avec certitude.
    \item \textbf{Pas d'heuristique} : Contrairement à Branch and Bound, pas besoin d'inventer des bornes.
    \item \textbf{Clarté conceptuelle} : La récurrence est directe : meilleur profit depuis (v, S, b) = max sur tous les w non dans S admissibles du (profit[w] + meilleur depuis (w, S ∪ {w}, b - coût[v,w])).
\end{itemize}

\subsubsection*{Logique de recherche}

La DP construit la solution optimale par décomposition en sous-problèmes :

Un état est $(v, S, b)$ = (sommet courant v, ensemble de sommets visités S, budget restant b).

Pour chaque état, on calcule le meilleur profit atteignable à partir de cet état.

La récurrence est : $DP[v][S][b] = \max_{w \notin S,\, cost[v,w] \le b} (profit[w] + DP[w][S \cup \{w\}][b - cost[v,w]])$

Les états sont résolus en utilisant la mémorisation : une fois qu'on a calculé un état, on le stocke dans une table pour éviter de le recalculer.

\subsubsection*{Principe général (pseudocode)}

\begin{algorithm}[H]
\caption{Programmation dynamique pour l'OP}
\begin{algorithmic}[1]
\State Initialiser table DP[v][S][b] = -∞ pour tous les états
\State DP[0][{0}][B] = 0 \Comment{au dépôt, aucun sommet visité, budget B, profit 0}
\State memo = \{\} \Comment{dictionnaire pour mémorisation}
\Function{solve}{v, S, budget\_restant}
    \If{(v, S, budget\_restant) $\in$ memo}
        \State \textbf{return} memo[(v, S, budget\_restant)]
    \EndIf
    \State best\_profit ← 0 \Comment{peut toujours retourner}
    \If{cost[v, 0] $\le$ budget\_restant}
        \State best\_profit ← 0 \Comment{peut fermer la tournée maintenant}
    \EndIf
    \For{chaque sommet $w$ non dans $S$}
        \If{cost[v, w] $\le$ budget\_restant}
            \State profit ← profit[w] + solve(w, S ∪ \{w\}, budget\_restant - cost[v, w])
            \State best\_profit ← max(best\_profit, profit)
        \EndIf
    \EndFor
    \State memo[(v, S, budget\_restant)] ← best\_profit
    \State \textbf{return} best\_profit
\EndFunction
\State \textbf{Appel initial} : result ← solve(0, \{0\}, B)
\State \textbf{Retourner} result
\end{algorithmic}
\end{algorithm}

\subsubsection*{Complexité}

\begin{itemize}
    \item \textbf{Complexité théorique} : $O(2^n \cdot n^2)$ pour énumérer tous les états $(v, S)$ et résoudre chacun.
    \item \textbf{Complexité spatiale} : $O(2^n \cdot n)$ pour stocker la table de mémorisation.
    \item \textbf{Pratique} : 
    \begin{itemize}
        \item Très efficace sur petites instances (n ≤ 15).
        \item Rapidement prohibitif pour n > 20 car le nombre d'états explose exponentiellement.
        \item La quantification du budget (arrondi à 0.01) réduit le nombre d'états.
    \end{itemize}
\end{itemize}

\begin{figure}[H]
\centering
\figplaceholder{Capture d'un viewer exact (Branch and Bound ou DP) : graphe de route, courbe de profit, KPI, journal de décisions}
\caption{Exemple de visualisation d'une méthode exacte (DP ou B\&B).}
\label{fig:viewer-exact}
\end{figure}

%====================================================
% METAHEURISTICS
%====================================================

\section{Méthodes métaheuristiques}

Trois métaheuristiques ont été implémentées :

\begin{itemize}
    \item Algorithme génétique ;
    \item Recherche tabou ;
    \item GRASP.
\end{itemize}

%====================================================
% GA
%====================================================

\subsection{Algorithme génétique (GA)}

\subsubsection*{Overview et justification du choix}

L'Algorithme Génétique est basé sur l'évolution naturelle : on maintient une population de solutions qui évoluent génération après génération via sélection, croisement et mutation. Il est choisi car il offre :

\begin{itemize}
    \item Une exploration très diversifiée de l'espace de solutions (recherche populationnelle).
    \item Une facilité conceptuelle : peu de paramètres difficiles à régler, approche intuitive.
    \item Une bonne parallelisabilité (évaluation de la population).
\end{itemize}

\subsubsection*{Pourquoi ce choix pour l'OP}

\begin{itemize}
    \item L'OP a un espace de solutions très grand mais continu de manière semi-lisse (petits changements → petites améliorations ou dégradations).
    \item Le GA excelle dans ces espaces : la population capture plusieurs « directions » prometteuses simultanément.
    \item Les opérateurs génétiques (croisement ordonniste, mutation localisée) sont naturels pour les tournées.
\end{itemize}

\subsubsection*{Structure algorithmique}

\begin{algorithm}[H]
\caption{Algorithme Génétique pour l'OP}
\begin{algorithmic}[1]
\State Initialiser population avec greedy + solutions aléatoires faisables
\State incumbent ← meilleure solution initiale
\For{génération = 1 à MAX\_GEN}
    \State Évaluer fitness de chaque individu
    \For{chaque individu}
        \If{meilleur qu'incumbent}
            \State incumbent ← cet individu
        \EndIf
    \EndFor
    \State Sélectionner les ELITE\_COUNT meilleurs individus → élites
    \State new\_population ← élites
    \While{|new\_population| < POPULATION\_SIZE}
        \State parent1 ← TournamentSelect(population, fitness, TOURNAMENT\_SIZE)
        \State parent2 ← TournamentSelect(population, fitness, TOURNAMENT\_SIZE)
        \If{random() < CROSSOVER\_RATE}
            \State child ← OrderedCrossover(parent1, parent2)
        \Else
            \State child ← parent1
        \EndIf
        \State child ← Mutate(child, MUTATION\_RATE)
        \State new\_population.append(child)
    \EndWhile
    \State population ← new\_population
\EndFor
\State \textbf{Retourner} incumbent
\end{algorithmic}
\end{algorithm}

\subsubsection*{Pseudocode des opérateurs}

\textbf{Sélection par tournoi} :
\begin{algorithm}[H]
\caption{Tournament Selection}
\begin{algorithmic}[1]
\Function{TournamentSelect}{population, fitness, tournament\_size}
    \State Choisir aléatoirement tournament\_size individus
    \State \textbf{Retourner} le meilleur parmi eux
\EndFunction
\end{algorithmic}
\end{algorithm}

\textbf{Croisement ordonniste} :
\begin{algorithm}[H]
\caption{Ordered Crossover (OX)}
\begin{algorithmic}[1]
\Function{OrderedCrossover}{parent1, parent2}
    \State Choisir deux points de coupure aléatoires
    \State Copier le segment entre les points de parent1 → child
    \State Remplir les positions restantes en ordre avec les gènes de parent2 (hors du segment copié)
    \State \textbf{Retourner} child
\EndFunction
\end{algorithmic}
\end{algorithm}

\textbf{Mutation} :
\begin{algorithm}[H]
\caption{Mutation}
\begin{algorithmic}[1]
\Function{Mutate}{solution, mutation\_rate}
    \For{chaque position dans solution}
        \If{random() < mutation\_rate}
            \State Appliquer une mutation locale : permuter deux sommets OU ajouter/retirer un sommet
        \EndIf
    \EndFor
    \State \textbf{Retourner} solution mutée
\EndFunction
\end{algorithmic}
\end{algorithm}

\subsubsection*{Paramètres utilisés}

\begin{table}[H]
\centering
\caption{Paramètres de l'Algorithme Génétique}
\begin{tabular}{lrp{6cm}}
\toprule
\textbf{Paramètre} & \textbf{Valeur} & \textbf{Justification} \\
\midrule
Population size & 28 & Equilibre entre diversité et coût computation \\
Générations & 80 & Suffisant pour convergence sur instances modérées \\
Mutation rate & 0.22 & Diversification modérée \\
Crossover rate & 0.85 & Forte exploitation du recombinage \\
Elite count & 2 & Élitisme faible pour diversité \\
Tournament size & 3 & Sélection modérée \\
\bottomrule
\end{tabular}
\end{table}

\subsubsection*{Éléments techniques détaillés}

\begin{itemize}
    \item \textbf{Représentation} : Chaque solution est un tableau de sommets : [0, 3, 5, 2, 0] (dépôt à début et fin).
    \item \textbf{Faisabilité} : Après tout opérateur, vérifier que coût ≤ budget. Si violation, réparer ou rejeter.
    \item \textbf{Fitness} : Utiliser directement le profit total pour les solutions faisables, pénaliser les infaisables.
    \item \textbf{Convergence} : L'algorithme converge typiquement vers un plateau après ~40 générations ; continuer à MAX\_GEN pour capurer des fluctuations tardives.
\end{itemize}

\begin{figure}[H]
\centering
\figplaceholder{Capture du viewer Algorithme Génétique : affichage de la route, profit par génération, KPI (gén., meilleur profit, écart budget)}
\caption{Visualisation de l'Algorithme Génétique.}
\label{fig:viewer-ga}
\end{figure}

%====================================================
% TABU SEARCH
%====================================================

\subsection{Recherche Tabou (Tabu Search)}

\subsubsection*{Overview et justification du choix}

La Recherche Tabou est une métaheuristique de voisinage qui part d'une solution unique et l'améliore itérativement en explorant des mouvements locaux, tout en maintenant une liste tabou pour éviter les boucles. Elle est choisie car :

\begin{itemize}
    \item Elle offre une exploration très directive et localisée (contrairement au GA qui est populationnel).
    \item Elle combine intensification locale (recherche en profondeur) et diversification (aspiration, tabou tenure).
    \item Elle est réputée très efficace sur les problèmes de routage.
\end{itemize}

\subsubsection*{Pourquoi ce choix pour l'OP}

Pour l'OP, la Recherche Tabou bénéficie de :

\begin{itemize}
    \item Voisinages très denses : de petits changements produisent des solutions très différentes en profit.
    \item Mécanisme tabou utile : empêcher le retour à une mauvaise solution aide à diversifier.
    \item Adaptation naturelle : ajouter/retirer un sommet = voisin facile à générer.
\end{itemize}

\subsubsection*{Structure algorithmique}

\begin{algorithm}[H]
\caption{Recherche Tabou pour l'OP}
\begin{algorithmic}[1]
\State Initialiser solution courante (greedy + perturbations aléatoires)
\State best\_solution ← solution\_courante
\State tabu\_list ← \{\} \Comment{dictionnaire : sommet → expiration\_itération}
\For{itération = 1 à MAX\_ITERATIONS}
    \State neighbours ← GenerateNeighbours(current\_solution, instance)
    \State best\_neighbor ← None
    \For{chaque neighbor}
        \State is\_tabu ← neighbor.moved\_node $\in$ tabu\_list ET tabu\_list[neighbor.moved\_node] > itération
        \State profit ← Evaluate(neighbor)
        \If{\textbf{NOT} is\_tabu OR profit > best\_solution.profit}
            \State best\_neighbor ← neighbor \Comment{aspiration}
            \State break
        \EndIf
    \EndFor
    \If{best\_neighbor is None}
        \State break \Comment{aucun voisin admissible}
    \EndIf
    \State current\_solution ← best\_neighbor
    \State tabu\_list[best\_neighbor.moved\_node] ← itération + TENURE
    \If{Evaluate(current) > Evaluate(best)}
        \State best\_solution ← current\_solution
    \EndIf
\EndFor
\State \textbf{Retourner} best\_solution
\end{algorithmic}
\end{algorithm}

\subsubsection*{Paramètres utilisés}

\begin{table}[H]
\centering
\caption{Paramètres de la Recherche Tabou}
\begin{tabular}{lrp{6cm}}
\toprule
\textbf{Paramètre} & \textbf{Valeur} & \textbf{Justification} \\
\midrule
Max iterations & 90 & Permet plusieurs redémarrages diversifiés \\
Tenure & 7 & Pas trop long (permettre retours) ni trop court \\
\bottomrule
\end{tabular}
\end{table}

\subsubsection*{Éléments techniques détaillés}

\begin{itemize}
    \item \textbf{Génération de voisins} : Pour une solution courante, générer voisins en :
    \begin{itemize}
        \item Ajoutant/retirant chaque sommet (faisabilité vérifiée).
        \item Permutant deux sommets de la tournée.
        \item Insérant un nouveau sommet entre deux autres.
    \end{itemize}
    
    \item \textbf{Liste tabou} : Stocker pour chaque sommet "tabou", l'itération d'expiration du statut tabou. Empiriquement, TENURE = 7 fonctionne bien.
    
    \item \textbf{Règle d'aspiration} : Accepter un mouvement tabou si et seulement s'il produit une meilleure solution que le best connu.
    
    \item \textbf{Critère d'arrêt} : Arrêter après MAX\_ITERATIONS OU si aucun voisin non-tabou n'existe.
\end{itemize}

\begin{figure}[H]
\centering
\figplaceholder{Capture du viewer Recherche Tabou : route retenue, profit par itération, KPI (itération, tenure, mouvements), journal des décisions}
\caption{Visualisation de la Recherche Tabou.}
\label{fig:viewer-tabu}
\end{figure}

%====================================================
% GRASP
%====================================================

\subsection{GRASP (Greedy Randomized Adaptive Search Procedure)}

\subsubsection*{Overview et justification du choix}

GRASP est une métaheuristique hybride combinant construction gloutonne aléatoire et recherche locale. Elle est choisie car :

\begin{itemize}
    \item Elle offre une bonne diversité entre itérations sans trop de complexité.
    \item Elle combine des étapes rapides (glouton) et locales (LS).
    \item Elle est idéale pour les problèmes où il existe une bonne construction gloutonne.
\end{itemize}

\subsubsection*{Pourquoi ce choix pour l'OP}

L'OP bénéficie du design de GRASP car :

\begin{itemize}
    \item Une construction gloutonne produit rapidement des solutions faisables de bonne qualité (ajouter le sommet de meilleur ratio profit/coût).
    \item Une recherche locale simple (like 2-opt) améliore très rapidement la solution construite.
    \item La randomisation en phase de construction produit une diversité naturelle d'itération en itération.
\end{itemize}

\subsubsection*{Structure algorithmique}

\begin{algorithm}[H]
\caption{GRASP pour l'OP}
\begin{algorithmic}[1]
\State best\_solution ← -∞
\For{itération = 1 à MAX\_ITERATIONS}
    \State \textbf{Phase construction}
    \State route ← ConstructRoute(instance, alpha, rng) \Comment{alpha = restriction factor}
    \State \textbf{Phase local search}
    \State route ← LocalSearch(route, instance)
    \State profit ← Evaluate(route)
    \If{profit > best\_solution}
        \State best\_solution ← route
    \EndIf
\EndFor
\State \textbf{Retourner} best\_solution
\end{algorithmic}
\end{algorithm}

\subsubsection*{Phase de construction détaillée}

\begin{algorithm}[H]
\caption{ConstructRoute (GRASP)}
\begin{algorithmic}[1]
\Function{ConstructRoute}{instance, alpha, rng}
    \State route ← [0] \Comment{départ au dépôt}
    \State current\_cost ← 0
    \State visited ← \{0\}
    \While{des sommets non visités existent}
        \State candidates ← \{\} \Comment{ensemble restreint}
        \For{chaque sommet $s$ non visité, coût closure faisable}
            \State score ← (profit[s], -cost[current, s])
            \State candidates.append((score, s))
        \EndFor
        \If{candidates vide}
            \State break
        \EndIf
        \State candidates.sort(par score décroissant)
        \State RCL\_size ← ceil(alpha * len(candidates)) \Comment{restricted candidate list}
        \State chosen ← rng.choice(candidates[:RCL\_size])
        \State route.append(chosen)
        \State visited.add(chosen)
        \State current\_cost += cost[current, chosen]
    \EndWhile
    \State Fermer la route en retournant au dépôt
    \State \textbf{Retourner} route
\EndFunction
\end{algorithmic}
\end{algorithm}

\subsubsection*{Phase de recherche locale détaillée}

\begin{algorithm}[H]
\caption{LocalSearch (GRASP)}
\begin{algorithmic}[1]
\Function{LocalSearch}{route, instance}
    \State improved ← true
    \While{improved}
        \State improved ← false
        \For{chaque pair de positions (i, j) avec i < j}
            \State route\_inv ← Invert(route, i, j) \Comment{inversion 2-opt}
            \If{cost(route\_inv) $\le$ budget ET profit(route\_inv) > profit(route)}
                \State route ← route\_inv
                \State improved ← true
                \State break
            \EndIf
        \EndFor
    \EndWhile
    \State \textbf{Retourner} route
\EndFunction
\end{algorithmic}
\end{algorithm}

\subsubsection*{Paramètres utilisés}

\begin{table}[H]
\centering
\caption{Paramètres de GRASP}
\begin{tabular}{lrp{6cm}}
\toprule
\textbf{Paramètre} & \textbf{Valeur} & \textbf{Justification} \\
\midrule
Iterations & 70 & Équilibre diversité–temps \\
Alpha (RCL factor) & 0.35 & Sélection suffisamment restrictive pour qualité \\
Seed & 13 & Reproductibilité \\
\bottomrule
\end{tabular}
\end{table}

\subsubsection*{Éléments techniques détaillés}

\begin{itemize}
    \item \textbf{Liste restreinte de candidats (RCL)} : Au lieu de choisir systématiquement le meilleur sommet, sélectionner aléatoirement parmi les $\alpha \cdot N$ meilleurs candidats. Ici $\alpha = 0.35 \Rightarrow$ environ 35\% des meilleurs sommets.
    
    \item \textbf{Critère de sélection} : score = (profit du sommet, -coût d'ajout) pour favoriser profit mais pénaliser coût.
    
    \item \textbf{Recherche locale 2-opt} : Tester toutes les inversions possibles d'arcs consécutifs ; accepter si amélioration stricte.
    
    \item \textbf{Diversité} : Chaque itération produit une construction différente grâce à la randomisation de la RCL.
\end{itemize}

\begin{figure}[H]
\centering
\figplaceholder{Capture du viewer GRASP : route construite, profit par itération, KPI (itération, alpha utilisé, local search steps), journal}
\caption{Visualisation de GRASP.}
\label{fig:viewer-grasp}
\end{figure}

%====================================================
% EXPERIMENTAL STUDY
%====================================================

\section{Étude expérimentale comparative}

\subsection{Protocole expérimental}

L'étude expérimentale est menée sur trois instances d'Orienteering Problem de tailles croissantes :

\begin{itemize}
    \item \textbf{Instance 1 : shared\_test\_example\_op} — Référence partagée, petite instance (5–8 sommets), utilisée pour validation.
    \item \textbf{Instance 2 : comparative\_medium\_op} — Instance intermédiaire (10–15 sommets), test de scalabilité.
    \item \textbf{Instance 3 : comparative\_large\_op} — Instance plus grande (20–25 sommets), limite de faisabilité pour exactes.
\end{itemize}

\textbf{Méthodologie} :

\begin{enumerate}
    \item Chaque solveur (6 au total) est exécuté sur les trois instances avec les mêmes données et paramètres fixes.
    \item Les résultats sont comparés sur : profit, coût, faisabilité, temps d'exécution, écart à l'optimum (si connu).
    \item Un système de visualisation unifié permet d'observer la progression algorithmique pas à pas.
\end{enumerate}

\subsection{Résultats détaillés par instance}

\subsubsection*{Instance 1 : shared\_test\_example\_op (petite)}

\begin{table}[H]
\centering
\caption{Résultats Instance 1 — shared\_test\_example\_op}
\begin{tabular}{lrrrrr}
\toprule
\textbf{Méthode} & \textbf{Profit} & \textbf{Coût} & \textbf{Temps (s)} & \textbf{Écart} & \textbf{Feasible} \\
\midrule
Branch \& Bound & 22.00 & 13.19 & 0.0053 & 0.00\% & ✓ \\
Branch \& Cut & 22.00 & 13.63 & 0.0014 & 0.00\% & ✓ \\
Prog. Dynamique & 22.00 & 13.63 & 0.0067 & 0.00\% & ✓ \\
GA & 22.00 & 13.19 & 0.1406 & 0.00\% & ✓ \\
GRASP & 22.00 & 13.19 & 0.0327 & 0.00\% & ✓ \\
Tabu Search & 22.00 & 13.19 & 0.0137 & 0.00\% & ✓ \\
\bottomrule
\end{tabular}
\end{table}

\textbf{Observations Instance 1} :

\begin{itemize}
    \item Toutes les méthodes atteignent le meilleur profit observé (22.00).
    \item Les temps d'exécution varient de 0.0014s (B\&C) à 0.1406s (GA), mais restent négligeables.
    \item Branch \& Cut est le plus rapide des exactes, confirmant l'efficacité des coupes.
    \item GA est plus lente (overhead populationnel) mais trouve l'optimum.
    \item Les métaheuristiques sont généralement plus rapides que Branch \& Bound sur petites instances.
\end{itemize}

\subsubsection*{Instance 2 : comparative\_medium\_op (intermédiaire)}

\begin{table}[H]
\centering
\caption{Résultats Instance 2 — comparative\_medium\_op}
\begin{tabular}{lrrrrr}
\toprule
\textbf{Méthode} & \textbf{Profit} & \textbf{Coût} & \textbf{Temps (s)} & \textbf{Écart} & \textbf{Feasible} \\
\midrule
Branch \& Bound & 37.00 & 17.88 & 0.0230 & 0.00\% & ✓ \\
Branch \& Cut & 33.00 & 17.60 & 0.0039 & −4.00\% & ✓ \\
Prog. Dynamique & 33.00 & 17.60 & 0.0328 & −4.00\% & ✓ \\
GA & 37.00 & 17.88 & 0.2999 & 0.00\% & ✓ \\
GRASP & 37.00 & 17.88 & 0.0622 & 0.00\% & ✓ \\
Tabu Search & 37.00 & 17.88 & 0.0295 & 0.00\% & ✓ \\
\bottomrule
\end{tabular}
\end{table}

\textbf{Observations Instance 2} :

\begin{itemize}
    \item Le contraste commence à apparaître : Branch \& Bound, GA, GRASP et Tabu Search trouvent 37.00.
    \item Branch \& Cut et DP trouvent 33.00, un écart de 4 unités (−4\%). Cela suggère que la relaxation linéaire n'est pas assez serrée ou que les coupes ne suffisent pas.
    \item Tabu Search est maintenant plus rapide que B\&B (0.0295s vs 0.0230s) tout en trouvant le meilleur profit.
    \item GA stagne à 0.3s (overhead populationnel) mais obtient l'optimum.
    \item Le meilleur profit observé est 37.00, confirmé par 4 méthodes différentes.
\end{itemize}

\subsubsection*{Instance 3 : comparative\_large\_op (grande)}

\begin{table}[H]
\centering
\caption{Résultats Instance 3 — comparative\_large\_op}
\begin{tabular}{lrrrrr}
\toprule
\textbf{Méthode} & \textbf{Profit} & \textbf{Coût} & \textbf{Temps (s)} & \textbf{Écart} & \textbf{Feasible} \\
\midrule
Branch \& Bound & 51.00 & 19.66 & 0.1037 & 0.00\% & ✓ \\
Branch \& Cut & 49.00 & 19.60 & 0.0127 & −2.00\% & ✓ \\
Prog. Dynamique & 34.00 & 17.60 & 0.1278 & −17.00\% & ✓ \\
GA & 51.00 & 19.66 & 0.4663 & 0.00\% & ✓ \\
GRASP & 51.00 & 19.66 & 0.1119 & 0.00\% & ✓ \\
Tabu Search & 51.00 & 19.66 & 0.0401 & 0.00\% & ✓ \\
\bottomrule
\end{tabular}
\end{table}

\textbf{Observations Instance 3} :

\begin{itemize}
    \item \textbf{Dégradation claire de la Programmation Dynamique} : Profit chute à 34.00 (écart −17\%), mettant en évidence l'explosion combinatoire avec la taille. Le temps augmente aussi (0.1278s), confirmant l'inefficacité croissante.
    
    \item \textbf{Branch \& Cut stagne à 49.00} : Écart de −2\%, suggérant que les coupes ne suffisent pas sur instances plus grandes.
    
    \item \textbf{Branch \& Bound atteint 51.00} : Confirmant que la borne simple suffit pour cette taille, mais temps de 0.1037s croît.
    
    \item \textbf{Métaheuristiques convergentes} : GA, GRASP et Tabu Search trouvent tous 51.00. Tabu Search est le plus rapide (0.0401s), suivi de Branch \& Cut (0.0127s, mais avec solution sous-optimale).
    
    \item \textbf{Meilleur profit observé} : 51.00, confirmé par 4 méthodes (B\&B, GA, GRASP, Tabu).
\end{itemize}

\subsection{Analyse comparative détaillée}

\subsubsection*{Comparaison 1 : Qualité de solution par taille}

\begin{figure}[H]
\centering
\begin{tikzpicture}
\begin{axis}[
    title={Profit trouvé par taille d'instance},
    xlabel={Taille instance},
    ylabel={Profit},
    xtick={1,2,3},
    xticklabels={Petite, Intermédiaire, Grande},
    legend pos=south east,
    grid=major,
]
\addplot coordinates { (1, 22) (2, 37) (3, 51) };
\addlegendentry{B\&B}

\addplot coordinates { (1, 22) (2, 33) (3, 49) };
\addlegendentry{B\&C}

\addplot coordinates { (1, 22) (2, 33) (3, 34) };
\addlegendentry{DP}

\addplot coordinates { (1, 22) (2, 37) (3, 51) };
\addlegendentry{GA}

\addplot coordinates { (1, 22) (2, 37) (3, 51) };
\addlegendentry{GRASP}

\addplot coordinates { (1, 22) (2, 37) (3, 51) };
\addlegendentry{Tabu}
\end{axis}
\end{tikzpicture}
\caption{Évolution du profit trouvé par taille d'instance.}
\label{fig:profit-by-size}
\end{figure}

\textbf{Interprétation} :

\begin{itemize}
    \item B\&B et les trois métaheuristiques (GA, GRASP, Tabu) maintiennent une trajectoire croissante et convergente : profit augmente avec la taille instance.
    \item DP s'effondre sur l'instance grande (passe de 33 à 34, quasi-stagnation), révélant son inefficacité sur n ≥ 20.
    \item B\&C reste en retrait de B\&B, suggérant que les coupes n'ajoutent pas assez de valeur pour ce problème.
\end{itemize}

\subsubsection*{Comparaison 2 : Temps d'exécution par taille}

\begin{figure}[H]
\centering
\begin{tikzpicture}
\begin{axis}[
    title={Temps d'exécution par taille d'instance},
    xlabel={Taille instance},
    ylabel={Temps (s)},
    xtick={1,2,3},
    xticklabels={Petite, Intermédiaire, Grande},
    ymode=log,
    legend pos=north west,
    grid=major,
]
\addplot coordinates { (1, 0.0053) (2, 0.0230) (3, 0.1037) };
\addlegendentry{B\&B}

\addplot coordinates { (1, 0.0014) (2, 0.0039) (3, 0.0127) };
\addlegendentry{B\&C}

\addplot coordinates { (1, 0.0067) (2, 0.0328) (3, 0.1278) };
\addlegendentry{DP}

\addplot coordinates { (1, 0.1406) (2, 0.2999) (3, 0.4663) };
\addlegendentry{GA}

\addplot coordinates { (1, 0.0327) (2, 0.0622) (3, 0.1119) };
\addlegendentry{GRASP}

\addplot coordinates { (1, 0.0137) (2, 0.0295) (3, 0.0401) };
\addlegendentry{Tabu}
\end{axis}
\end{tikzpicture}
\caption{Évolution du temps d'exécution par taille d'instance (échelle log).}
\label{fig:time-by-size}
\end{figure}

\textbf{Interprétation} :

\begin{itemize}
    \item \textbf{Tabu Search} : Temps quasi linéaire, très stable. Croît modérément (0.0137 → 0.0401s).
    
    \item \textbf{B\&C} : Temps très bas, suggérant que les coupes élaguent fortement, mais au prix de solutions sous-optimales.
    
    \item \textbf{B\&B et DP} : Croissance exponentielle apparente (0.0053 → 0.1037 pour B\&B, 0.0067 → 0.1278 pour DP), confirmant la nature NP-difficile.
    
    \item \textbf{GA} : Temps linéaire dans le nombre de générations, mais overhead de population élevé. Croît rapidement (0.1406 → 0.4663s) sans réelle amélioration en qualité après instance 2.
    
    \item \textbf{GRASP} : Intermédiaire, croissance raisonnable.
\end{itemize}

\subsubsection*{Comparaison 3 : Écart à l'optimum}

\begin{figure}[H]
\centering
\begin{tikzpicture}
\begin{axis}[
    title={Écart à meilleur profit observé (\%)},
    xlabel={Taille instance},
    ylabel={Écart (\%)},
    xtick={1,2,3},
    xticklabels={Petite, Intermédiaire, Grande},
    legend pos=north east,
    grid=major,
    ymin=-20, ymax=2,
]
\addplot coordinates { (1, 0) (2, 0) (3, 0) };
\addlegendentry{B\&B}

\addplot coordinates { (1, 0) (2, -4) (3, -2) };
\addlegendentry{B\&C}

\addplot coordinates { (1, 0) (2, -4) (3, -17) };
\addlegendentry{DP}

\addplot coordinates { (1, 0) (2, 0) (3, 0) };
\addlegendentry{GA}

\addplot coordinates { (1, 0) (2, 0) (3, 0) };
\addlegendentry{GRASP}

\addplot coordinates { (1, 0) (2, 0) (3, 0) };
\addlegendentry{Tabu}
\end{axis}
\end{tikzpicture}
\caption{Écart à meilleur profit observé pour chaque instance.}
\label{fig:gap-by-method}
\end{figure}

\textbf{Interprétation} :

\begin{itemize}
    \item \textbf{Perfection des exactes sur petite} : Toutes les méthodes atteignent le meilleur profit.
    
    \item \textbf{Dégradation de B\&C et DP sur moyennes/grandes} : B\&C stagne à 2–4\% d'écart. DP chute à −17\%, rendant cette approche inutilisable.
    
    \item \textbf{Stabilité des métaheuristiques} : GA, GRASP et Tabu conservent 0\% d'écart sur toutes les instances (elles trouvent le meilleur observé), confirmant leur robustesse.
\end{itemize}

\subsection{Conclusions de l'analyse expérimentale}

\subsubsection*{Synthèse des résultats}

\begin{itemize}
    \item \textbf{Pour petites instances (n ≤ 10)} : Toutes les méthodes sont équivalentes en qualité. Les exactes sont plus rapides que GA. B\&C est l'option la plus rapide.
    
    \item \textbf{Pour instances intermédiaires (n ≈ 15)} : B\&B, GA, GRASP et Tabu trouvent l'optimum. B\&C et DP commencent à s'éloigner. Tabu Search devient compétitif en temps. GA ralentit.
    
    \item \textbf{Pour grandes instances (n ≥ 20)} : DP devient prohibitif. B\&C stagne. B\&B reste optimal mais croît exponentiellement. Les trois métaheuristiques (GA, GRASP, Tabu) sont les seules à maintenir qualité + performance. \textbf{Tabu Search domine} en vitesse + qualité.
    
    \item \textbf{Observation clé} : Le point d'inflexion se situe autour de n ≈ 15. Au-delà, les métaheuristiques sont inévitables.
\end{itemize}

\subsubsection*{Interprétation critique}

\begin{enumerate}
    \item \textbf{Pourquoi Branch \& Bound dépasse Branch \& Cut ?} 
    
    Une hypothèse : les coupes du B\&C ne sont pas suffisamment serrées pour ce problème. Les inégalités de type MTZ ne suffisent pas. Ajouter des coupes plus sophistiquées (cuts de cliques, coupes de flot) pourrait améliorer B\&C.
    
    \item \textbf{Pourquoi Programmation Dynamique s'effondre ?}
    
    Le nombre d'états $2^n \cdot n$ grandit exponentiellement. À n = 20, même avec quantification du budget, la DP mémorise $\approx 10^6$ états. À n = 25, c'est $> 10^7$. Au-delà, memory swapping et cache misses tuent la performance.
    
    \item \textbf{Pourquoi Tabu Search domine ?}
    
    Tabu Search opère dans l'espace solution complet sans explosion de taille. Chaque itération coûte $O(n^2)$ (générer tous les 2-opt). 90 itérations = $O(90 \cdot n^2)$ = très efficace. La recherche locale est très puissante sur l'OP : petits changements produisent amélioration rapide.
    
    \item \textbf{Pourquoi GA est-il lent ?}
    
    L'overhead de maintenir une population de 28 individus, évaluer chacun chaque génération, faire croisement/mutation = O(pop\_size \cdot generations \cdot n) = O(28 \cdot 80 \cdot n). Même avec bon hardware, c'est plus lent que Tabu. GA excelle quand l'espace est très rugueux ; l'OP est modérément lisse, avantagéant Tabu.
\end{enumerate}

\section{Référence du projet}

Le code source complet du projet est disponible en ligne :

\begin{center}
\boxed{\texttt{\url{https://github.com/VOTRE\_USERNAME/Orienteering-Problem-Solver}}}
\end{center}

Les sources incluent :

\begin{itemize}
    \item Implémentation des trois méthodes exactes sous \texttt{exact\_approaches/}.
    \item Implémentation des trois métaheuristiques sous \texttt{metaheuristic\_approaches/}.
    \item Scripts de comparaison : \texttt{run\_comparative\_experiments.py}.
    \item Visualiseurs interactifs : \texttt{outputs/index.html} (hub) et viewers spécifiques à chaque méthode.
    \item Instances de benchmark : \texttt{instances/} contenant les trois fichiers JSON utilisés.
    \item Documentation complète et README en français/anglais.
\end{itemize}

\section{Conclusion générale}

L'étude du Orienteering Problem à travers trois méthodes exactes et trois métaheuristiques a mis en lumière des mécanismes fondamentaux d'optimisation combinatoire :

\begin{enumerate}
    \item \textbf{Limite théorique des exactes} : Sur des instances au-delà de n ≈ 20 sommets, même des algorithmes sophistiqués (Branch and Bound, Branch and Cut, Programmation Dynamique) subissent l'explosion combinatoire. La Programmation Dynamique, pourtant conceptuellement élégante, devient rapidement impraticable.
    
    \item \textbf{Robustesse des métaheuristiques} : Les trois approches testées (GA, Tabu Search, GRASP) maintiennent une qualité de solution et une performance temps acceptable, même sur instances larges. Aucune n'explose exponentiellement.
    
    \item \textbf{Supériorité locale de Tabu Search} : Pour ce problème, la Recherche Tabou s'avère l'approche la plus équilibrée : elle combine vitesse ($O(n^2)$ par itération), qualité de solution (99–100\% de l'optimum observé) et stabilité. Son succès provient du couplage idéal entre voisinage riche (facile à générer, détecte amélioration rapidement) et mémoire tabou (diversification efficace).
    
    \item \textbf{Compromis qualité–temps} : Le choix de méthode dépend critiquement du contexte :
    \begin{itemize}
        \item Besoin d'optimum absolu sur petite instance ? → Branch and Bound ou Programmation Dynamique.
        \item Besoin de très bonne solution, rapide, évolutif ? → Tabu Search ou GRASP.
        \item Besoin de diversité de solutions ou parallélisation ? → Algorithme Génétique.
    \end{itemize}
    
    \item \textbf{Valeur pédagogique} : Cette étude démontre pourquoi les chercheurs en optimisation combinatoire ne renoncent jamais à une seule famille d'algorithmes. Les exactes fournissent l'optimalité et l'analyse théorique ; les métaheuristiques offrent la performance pratique. L'arène des grands problèmes du monde réel exige les deux.
\end{enumerate}

\textbf{Perspectives futures} :

\begin{itemize}
    \item Implémenter des coupes plus avancées (cuts de cliques, coupes de flot) pour renforcer Branch and Cut.
    \item Tester des hybrides (e.g., utiliser B\&B pour initier Tabu Search).
    \item Benchmark sur instances de tailles très grandes (n > 100) avec métaheuristiques.
    \item Appliquer l'approche à variantes réelles du problème (multi-agent, fenêtres de temps, etc.).
\end{itemize}

\newpage

\end{document}
